%% kapa: exam=1 round=27 year=2016 period=1 form=A subject=civil_law subject_ko=민법 seq=30 number=30
%% source: Q-Net 감정평가사 기출문제 (cst00309) artlSeq=5169402 file=extracted/1008/5169402/제27회 감정평가사 1차 문제/제27회 감정평가사 1교시문제(A형).pdf
%% conversion: kordoc 4.24.51 -> tools/gen1.py (auto); review=proofread-sonnet; flags=none
\begin{question}{30}
甲 소유 X토지에 대하여 乙이 점유취득시효를 완성하였으나 등기를 경료하지 못하고 있는 경우에 관한 설명으로 옳지 않은 것은? (다툼이 있으면 판례에 따름)
\begin{choices}
  \item 甲이 丙에게 X토지를 매도하여 이전등기를 마치면, 乙은 甲에 대한 시효취득의 효력을 丙에게 주장할 수 없다.
  \item 위의 ①에서 丙이 甲의 배임행위에 적극 가담한 경우에는 甲과 丙의 매매는 반사회질서 법률행위로서 무효가 된다.
  \item 乙이 점유를 상실하면 시효이익의 포기로 간주되어 취득한 소유권이전등기청구권은 소멸한다.
  \item 乙의 X토지에 대한 취득시효의 주장에도 불구하고 甲이 악의로 丙에게 이를 매도한 경우, 乙은 甲에 대하여 손해배상을 청구할 수 있다.
  \item X토지가 수용된 경우, 그 전에 乙이 甲에 대하여 시효취득기간만료를 원인으로 등기청구권을 행사하였다면 대상청구권을 행사할 수 있다.
\end{choices}
\end{question}
\answer{3}
